% !TeX program = pdflatex
% !TeX TXS-program:compile = txs:///pdflatex/[-shell-escape]
% !TeX TXS-program:pdflatex = pdflatex -synctex=1 -interaction=nonstopmode --shell-escape %.tex

\documentclass[../../../main_compilation.tex]{subfiles}

\begin{document}

% ================================================================================
% Encabezado / Ficha de la Sala TryHackMe
% ================================================================================
\section[THM - NOMBRE\_SALA (10.10.X.X)]{NOMBRE\_SALA — TryHackMe}
\label{sec:thm-nombre-sala}

\targetSummary%
{NOMBRE\_SALA}% Nombre de la sala
{10.10.X.X}% IP
{TryHackMe}% Plataforma
{Medium}% Dificultad: Easy, Medium, Hard, Insane
{Linux}% Sistema Operativo: Linux, Windows, Active Directory
{Web / PrivEsc / Forensics}% Categoría o Tags clave
{\today}% Fecha de resolución
{thm\{user\_flag\}}% User Flag
{thm\{root\_flag\}}% Root Flag

\vspace{0.4cm}

% ================================================================================
% 1. Descripción de la Sala
% ================================================================================
\subsection{Descripción del Reto}
Objetivos de la sala de TryHackMe y conceptos evaluados (ej. Active Directory, Web Application Security, Network Security, etc.).

\begin{vulnbox}{Vulnerabilidad / Vector Principal}{High}
	\textbf{Tipo:} Desconfiguración / Inyección / CVE-XXXX \hfill \textbf{Severidad:} High\\
	\textbf{Descripción:} Resumen del fallo explotado durante la sala.
\end{vulnbox}

% ================================================================================
% 2. Tareas y Pasos de Resolución
% ================================================================================
\subsection{Reconocimiento \& Tareas Iniciales}

\begin{terminal}[kali@pentest:\~{}]
	nmap -sC -sV -Pn 10.10.X.X -oN nmap_thm.txt
\end{terminal}

\subsection{Explotación y Acceso}

\begin{terminal}[kali@pentest:\~{}]
	# Comandos de explotación o payloads ejecutados
\end{terminal}

\subsection{Escalada de Privilegios}

\begin{terminal}[user@target:\~{}]
	# Procedimiento para obtener root / SYSTEM
\end{terminal}

\begin{flagbox}{User Flag}
	thm{user_flag_aqui}
\end{flagbox}

\begin{flagbox}{Root Flag}
	thm{root_flag_aqui}
\end{flagbox}

% ================================================================================
% 3. Conceptos Aprendidos y Mitigaciones
% ================================================================================
\subsection{Lecciones Aprendidas \& Mitigación}

\begin{mitigationbox}[Buenas Prácticas]
	\begin{itemize}[leftmargin=*]
		\item Medidas recomendadas para evitar esta vulnerabilidad en producción.
		\item Endurecimiento y principios de mínimo privilegio aplicables.
	\end{itemize}
\end{mitigationbox}

\begin{learningbox}[Conocimientos Técnicos Adquiridos]
	\begin{itemize}[leftmargin=*]
		\item \textbf{Concepto Clave:} Técnica o vector principal practicado en la sala.
		\item \textbf{Táctica / Herramienta:} Uso de utilidades específicas o comandos de enumeración.
	\end{itemize}
\end{learningbox}

\begin{sourcebox}[Fuentes de Referencia \& Documentación]
	\begin{itemize}[leftmargin=*]
		\resourceItem{Docs}{Documentación de la Sala (THM)}{Enlace o manual de la sala}
		\resourceItem{Web}{Recurso / Cheatsheet de Apoyo}{\url{https://...}}
	\end{itemize}
\end{sourcebox}

\end{document}
